\subsection{The hypothetical center is simple}
\label{subsec:r3472-focus-simple-center}

Write $R=\sha{R8_3_3472}$ and $r=\sqrt2$. Its positive dimension
character, in canonical order, is
\[
 (1,6+5r,6+5r,7+5r,6+5r,8+5r,8+5r,8+5r),
 \qquad D=700+490r.
\]
An exhaustive check of the $128$ unit-containing subsets shows that
the only based fusion subrings are $\mathbb Z\one$ and $R$.
Independently, every nonunit fusion matrix has entrywise-positive square.
The regular Casimir has characteristic polynomial
\[
 (x-7)^6(x^2-1400x+9800).
\]
Consequently, in a pseudo-unitary categorification, the distinguished
summands of $I(\one)$ have dimensions
\[
 1,\quad 99+70r,\quad
 \underbrace{100+70r,\ldots,100+70r}_{6\text{ copies}}.
\]

\begin{proposition}
If $R$ has a complex categorification $\mathcal C$, then
$Z(\mathcal C)$ is simple. Every nonunit central simple has nonintegral
Frobenius--Perron dimension. In particular its center has no nontrivial
pointed, integral, or Tannakian fusion subcategory.
\end{proposition}
\begin{proof}
Use the previously established general sphericalization and Galois
reduction for $R$. If $\mathcal C$ were braided, its braiding would
transport to its pseudo-unitary Galois conjugate. In a braided category,
the squared norm of a simple is a $d$-number. In the pseudo-unitary
conjugate this is its squared Frobenius--Perron dimension, and an
algebraic integer is a $d$-number if and only if its square is one.
The dimension $6+5r$ has trace $12$ and norm $-14$; the quadratic
$d$-number divisibility condition fails because $14\nmid144$.
Thus no categorification of $R$ is braided.

Liu--Palcoux--Ren \cite[Theorem~5.1]{LPRSimpleFocus} state that a nonpointed simple
complex fusion category has simple center if and only if it admits
no braiding. The theorem does not assume integrality, despite the
title of the paper. The verified based simplicity of $R$ therefore
proves the first assertion.

For the second assertion let $Z$ be a central simple of integral
FP dimension, and write $F(Z)=\sum_i n_iX_i$. The coefficient of $r$
in its FP dimension is $5\sum_{i>0}n_i$. It follows that $F(Z)$ is
a sum of units. Every object in the fusion subcategory generated by
$Z$ consequently also forgets to a sum of units. If $Z\ne\one$,
simplicity would make this the entire center, contrary to the
dominance of the forgetful functor onto the nontrivial category
$\mathcal C$. Hence $Z=\one$.
\end{proof}

These assertions are restrictions on a hypothetical category, not an
exclusion of $R$. In particular the integral prime-power central-simple
obstruction in Remark~5.15 of the cited paper cannot be applied here:
its hypotheses require an integral modular category. The center here
has irrational global FP dimension $D^2$.

The exact ring checks are recorded in
\begin{quote}\small
\path{evidence/r3472_focus_20261022/center/verify_structure.py}.
\end{quote}
They recheck the tensor, based subrings, positive dimension character,
formal codegrees and primitive fusion matrices. They do not replace
the earlier complete signed-quotient proof of the spherical reduction.


\subsection{Exact rational and integer center witnesses}
\label{subsec:r3472-focus-rational-center}

The previously published strengthened center system for
\sha{R8_3_3472} has $3347$ equations and $13242$ variables. Its
constraints comprise the induction Gram matrix, distinguished
$S$-symmetry, cubic Verlinde equations, Galois spectral balance,
Frobenius--Perron bounds, and unit-product constraints. The certificate
\begin{quote}\small
\path{evidence/r3472_focus_20261022/center/rational_witness.json}
\end{quote}
gives an exact nonnegative rational solution satisfying every propagated
box bound. It has $234$ nonzero entries, all integral except one:
the high-column multiplicity for restriction vector
$(0,10,10,10,10,10,10,10)$ equals $49/50$.

The accompanying verifier reconstructs the unchanged published system
from the original census, checks its source inventory and matrix hashes,
replays the box-bound derivation, and verifies all equations and bounds
by rational arithmetic. There is also an exact integer solution to this same system. Let
$u=(0,1,1,1,1,1,1,1)$. The complete column for $10u$ equals $100$ times
the column for $u$. Replacing weight $49/50$ on $10u$ by weight $98$
on $u$ preserves all equations and satisfies all propagated bounds.
Every other coordinate was already integral. The certificate
\begin{quote}\small
\path{evidence/r3472_focus_20261022/center/integer_witness.json}
\end{quote}
records the resulting $234$ nonzero integer entries. The same verifier
checks the full column identity and all equations and bounds for both
witnesses in one reconstruction.

Thus this particular necessary relaxation cannot exclude $R$, even by
integer infeasibility. These certificates do not incorporate twist-phase
equations and do not provide a modular datum or categorification.
Categorifiability remains unresolved.
